% !TeX program = pdflatex
\documentclass[12pt]{exam}

\usepackage[margin=0.5in]{geometry}
\usepackage{listings}
\usepackage{tcolorbox}
\usepackage{xcolor}
\usepackage[T1]{fontenc}
\usepackage[utf8]{inputenc}
\newcommand{\emojifont}{}
\setlength{\parindent}{0pt}

\definecolor{codegray}{rgb}{0.5,0.5,0.5}
\definecolor{lightgray}{gray}{0.95}
\lstset{
  language=SQL,
  basicstyle=\ttfamily\small,
  keywordstyle=\color{blue},
  commentstyle=\color{codegray},
  backgroundcolor=\color{lightgray},
  frame=none,
  breaklines=true,
  showstringspaces=false
}

\begin{document}
\pagestyle{empty}

\textbf{Database Management Systems} \hfill Referential Integrity with Triggers -- Worksheet 1 \\
Beata Kubiak, Instructor

\medskip

Let's use triggers to ensure that each \texttt{time\_slot\_id} in the \texttt{section} table corresponds to a \texttt{time\_slot\_id} in the \texttt{time\_slot} table. The two triggers below were presented during the lecture and contribute to this goal. However, additional triggers are needed to fully enforce the constraint. Please implement them.

Submit your work in a file named \texttt{CSC3300\_referential-integrity-with-triggers.sql} and upload it to the Dropbox in iLearn under the title \textbf{Bonus: Referential Integrity with Triggers}.

\medskip

\textbf{Trigger: Before Inserting into Section}
\begin{tcolorbox}[colback=gray!5, colframe=black!60]
\begin{lstlisting}
create trigger timeslot_check1
before insert on section
for each row 
begin
    if (new.time_slot_id not in (
        select time_slot_id
        from time_slot))
    then
        SIGNAL sqlstate '45001' set message_text = "Cannot insert section: time_slot_id does not exist in time_slot.";
    end if;
end
\end{lstlisting}
\end{tcolorbox}

\medskip

\textbf{Trigger: After Deleting from Time\_Slot} % <-- underscore escaped
\begin{tcolorbox}[colback=gray!5, colframe=black!60]
\begin{lstlisting}
create trigger timeslot_check2
after delete on time_slot
for each row 
begin
    if (old.time_slot_id not in (
        select time_slot_id
        from time_slot) and 
        old.time_slot_id in (  
        select time_slot_id
        from section))
    then
        SIGNAL sqlstate '45001' set message_text = 
        "Cannot delete this time slot: one or more course sections reference it.";
    end if;
end
\end{lstlisting}
\end{tcolorbox}

\end{document}
